\section{Problem Definition}\label{sec:problem-definition}

\subsection{Description}
% Describe stations, loading depots, products, vehicles, home garages, repeated
% loading/delivery operations, inputs, decisions, and planning horizon in plain
% language. State whether a load contains one product and is fully unloaded.

\subsection{Notation}
% Define common sets and parameters before model-specific variables. Keep units
% visible and use the same notation in the MILP and hybrid approach.

\begin{table}[H]
  \centering
  \caption{Common sets and indices}
  \label{tab:common-sets}
  \begin{tabularx}{\textwidth}{@{}lX@{}}
    \toprule
    Symbol & Definition \\
    \midrule
    $S$ & Set of physical delivery stations, indexed by $s$ \\
    $D$ & Set of physical loading depots, indexed by $d$ \\
    $P$ & Set of products, indexed by $p$ \\
    $K$ & Set of vehicles, indexed by $k$ \\
    $R$ & Request set, $R=\{(s,p)\in S\times P:q_{sp}>0\}$ \\
    $U$ & Supply-node set, $U=\{(d,p)\in D\times P:S_{dp}>0\}$ \\
    \bottomrule
  \end{tabularx}
\end{table}

\begin{table}[H]
  \centering
  \caption{Common parameters}
  \label{tab:common-parameters}
  \begin{tabularx}{\textwidth}{@{}lXl@{}}
    \toprule
    Symbol & Definition & Unit \\
    \midrule
    $q_{sp}$ & Demand for product $p$ at station $s$ & \emph{Specify} \\
    $S_{dp}$ & Available stock of product $p$ at depot $d$ & \emph{Specify} \\
    $Q_k$ & Capacity of vehicle $k$ & \emph{Specify} \\
    $c_{ij}$ & Travel cost between physical locations $i$ and $j$ & cost unit \\
    $\gamma_{pp'}$ & Cost of changing from product $p$ to product $p'$ & cost unit \\
    \bottomrule
  \end{tabularx}
\end{table}

\subsection{Requests and Supply Nodes}
% Explain that logical nodes retain their physical location and product, that
% several logical nodes may share a location, and that requests may be split.

\begin{description}
  \item[Station] A physical delivery location.
  \item[Depot] A physical loading location.
  \item[Request] A station--product pair with positive demand.
  \item[Supply node] A depot--product pair with available stock.
  \item[Delivery piece] A portion of a request created by demand splitting.
  \item[Mini-trip] One loading operation followed by deliveries of that product,
        constructed by the hybrid method.
\end{description}

\subsection{Assumptions and Constraints}
% Separate operating assumptions from enforced feasibility rules. Address:
% demand satisfaction, stock, capacity, product consistency, loading/delivery,
% continuity, departure/return, demand splitting, and every other actual rule.
% Identify any restriction that belongs only to the hybrid method.

\subsection{Objective Function}
% Define total travel plus changeover cost, including initial product states,
% same-product transitions, final cleaning, asymmetry, and any loading charge.
% Explain cost units and give one complete small-route calculation.
\begin{equation}
  \min C_{\mathrm{total}}
  = C_{\mathrm{travel}} + C_{\mathrm{changeover}}.
  \label{eq:common-objective}
\end{equation}

\section{MILP Model}\label{sec:milp-model}

\subsection{Decision Variables}
% Define every MILP variable and its domain. Use a table if the list is long.

\subsection{Objective Function}
% Expand Equation~\eqref{eq:common-objective} using the MILP variables.

\subsection{Constraints}
% Group equations by operational purpose and explain each family after it.
\subsubsection{Demand Satisfaction}
\subsubsection{Stock Availability}
\subsubsection{Vehicle Capacity and Product Consistency}
\subsubsection{Loading, Delivery, and Product-State Transitions}
\subsubsection{Route Continuity, Departure, and Return}
\subsubsection{Demand Splitting and Visit Limits}
% Justify finite visit/trip limits and every big-M value. Keep the complete core
% formulation here; move only lengthy derivations or strengthening details to an
% appendix.

\section{Proposed Approach}\label{sec:proposed-approach}

\subsection{Overview}
% Show: requests/supply nodes -> delivery pieces -> mini-trips -> CP assignment
% and sequencing. State exactly which decisions construction fixes and the
% resulting loss of flexibility.

\begin{table}[H]
  \centering
  \caption{Allocation of decisions between construction and constraint programming}
  \label{tab:decision-allocation}
  \begin{tabularx}{\textwidth}{@{}lXX@{}}
    \toprule
    Decision & Construction heuristic & CP model \\
    \midrule
    Request splitting & \emph{Specify} & \emph{Specify} \\
    Supply-node selection & \emph{Specify} & \emph{Specify} \\
    Deliveries within a mini-trip & \emph{Specify} & \emph{Specify} \\
    Vehicle assignment & \emph{Specify} & \emph{Specify} \\
    Mini-trip order & \emph{Specify} & \emph{Specify} \\
    \bottomrule
  \end{tabularx}
\end{table}

\subsection{Mini-Trip Construction}
% Specify inputs, outputs, ordering, splitting, stock updates, packing, internal
% routing, tie-breaking, guarantees, and failure behavior.
\begin{algorithm}[H]
  \KwData{Requests $R$, supply nodes $U$, fleet capacities, travel costs}
  \KwResult{A set of mini-trips, or a reported construction failure}
  \tcp{Replace this scaffold with verified construction pseudocode.}
  order requests according to the implemented rule\;
  split and pack requests while updating available stock\;
  select a supply node for each mini-trip\;
  order the deliveries within each mini-trip\;
  \Return{constructed mini-trips}\;
  \caption{Mini-trip construction scaffold}\label{alg:minitrip-construction}
\end{algorithm}

\subsection{CP Model}
% Define one sequence per vehicle, exactly-once mini-trip assignment, capacity
% compatibility, split-request restrictions, and all connector/internal costs.
% Explain which constraints construction has already enforced.

\subsection{Search Strategy}
% Separate feasibility propagation, search ordering, and cost bounding. Describe
% first-fail selection, minimum-detour placement, insert/forbid branching,
% incumbents, and stopping conditions as actually implemented.

\subsection{Implementation}
% Record MaxiCP and project versions. Document the original TransitionCost file
% and commit, retained license/attribution, port changes, and problem-specific
% transition matrix. Do not present the upstream constraint as a new invention.

\section{Experimental Setup}\label{sec:experimental-setup}

\subsection{Instances}
% Document the 100 synthetic instances: configurations, distributions, spatial
% design, demands, stocks, capacities, product states, changeover matrices,
% seeds, feasibility checks, and rejection/repair rules.

\subsection{Environment and Validation}
% Hardware, languages, exact solver/code versions and commits, threads, budgets,
% seeds, and independent checker rules. Include construction in hybrid runtime.

\subsection{Experiments and Metrics}
% Predefine the paired method comparison and the controlled with/without-
% changeover experiment. Define statuses, denominators, percentage signs, and
% zero handling. A reduced-CP bound is not automatically an original-problem
% bound.
